\minitopic{预测问题与干预问题不能共用一个答案}一个模型可以准确预测哪些患者会再次入院，却没有说明哪项措施能够降低再入院率。预测寻找变量与结果之间稳定的关联；干预判断则要比较“实施某措施”与“不实施”造成的差异。若贫困程度能预测健康风险，不表示改变数据库中的收入字段就会改善健康；若接受某治疗的人恢复较慢，也可能因为病情更重者更常接受治疗。把预测权重解释成因果机制，会从形式正确的模型推出错误行动。

决定是否干预时，应画出可能的共同原因、选择机制和中介路径，并寻找随机实验、自然实验或可靠机制证据。观察数据仍然有价值，但结论强度必须与设计相称。预测模型回答“接下来可能发生什么”，因果模型尝试回答“若我们改变某项因素会怎样”；解释模型还可能回答“哪些关系使结果可理解”。三类任务可以使用同一数据，却要求不同验证。

\minitopic{分布变化使旧校准失效}模型在去年、甲医院或某年龄组中校准良好，不保证换到今年、乙医院或另一群体仍然校准。疾病流行率改变会移动先验，设备更新会改变测量误差，使用者知道模型规则后还会调整行为。部署不是把一个固定工具放进不变世界，而是让模型进入会回应它的社会系统。模型一旦影响谁被检查、谁获得资源，下一轮数据就部分由旧模型制造。

因此，验证不能只在上线前完成。机构应预先确定哪些变化触发重新估计：基础率越过什么范围，缺失率或输入分布怎样漂移，哪类错误突然增加，政策目标是否改变。监测到差异后也不能盲目重训；新数据若已受旧决策污染，直接学习可能加深反馈循环。需要保留未受模型选择的审计样本、记录政策变动，并区分世界改变与记录方式改变。

\minitopic{稳健决定不必等待唯一最佳模型}面对多个都能解释现有数据的模型，一种反应是继续搜集证据，直到选出唯一赢家；现实决定常没有这么多时间。另一种方法是比较不同模型下各行动的后果。若“加固堤坝”在多种合理暴雨模型下都避免巨大损失，且成本可接受，它可能是稳健行动；若某措施只在一个脆弱模型中显得最优，就应优先搜集能区分模型的证据。

稳健不等于永远选择最保守方案。过度预防也有机会成本，反复误报会削弱未来信任。可以先列出合理模型集合，再检查每项行动的最坏损失、后悔值、可逆性和后续学习机会。可逆的小步试验有时优于一次性最大化期望效用，因为它同时行动并产生新信息。

\begin{argumentbox}
模型进入行动前要通过三道门：预测门检查新数据上的校准和失效群体；因果门检查建议的干预是否真能改变结果；治理门检查决定目标、权利限制、申诉和重新评估。通过第一道门的高分模型，仍可能在后两道门失败。
\end{argumentbox}

\index{预测模型}
\index{因果模型}
\index{分布变化}
\index{稳健决策}
